\documentstyle[%


righttag%


,amssymb%


,11pt%


,russian%


,izvru%


]{amsart}





%\setlength{\textwidth}{170mm}


%\setlength{\textheight}{232mm}


%\setlength{\hoffset}{-18mm}


%\setlength{\voffset}{-12mm}


%\setlength{\parindent}{6mm}





\newcommand{\tl}[3]{\noindent{\bf#1}\ #2 \dotfill #3 \par} 


\setcounter{page}{100}


\pagestyle{plain}


\begin{document}


\par


\vskip 2cm


\par


\bigskip


\bigskip


\bigskip


\par


{\centerline{\Large СОДЕРЖАНИЕ ЖУРНАЛА ``ИЗВЕСТИЯ ВЫСШИХ УЧЕБНЫХ ЗАВЕДЕНИЙ''}}





{\centerline{\Large серия ``Математика'' за 1997 г.}}


\bigskip


\bigskip


%%%%%%%%%% Use \newline %%%%%%%%%%%








\tl{Абдурагимов Э.И.} 


{О положительном радиально-симметрическом решении задачи\newline


Дирихле для одного нелинейного уравнения 


и численном методе его получения}{\No\,5}


\tl{Азбелев Н.В., Симонов П.М.}


{Устойчивость уравнений с запаздывающим аргументом}{\No\,6}


\tl{Амосов А.А., Злотник А.А.}{Полудискретный метод решения уравнений


одномерного\newline движения


неоднородного вязкого теплопроводного газа с негладкими данными}{\No\,4}


\tl{Антипин А.С.}{Вычисление неподвижных точек


симметричных экстремальных\newline отображений}{\No\,12}


\tl{Апекина Е.В., Хамисов О.В.}{Модифицированный метод симплексных


погружений\newline


с одновременным введением нескольких секущих плоскостей}{\No\,12}


\tl{Баскаков А.Г.}{Спектральный анализ разностных операторов 


с сильно непрерывными \newline коэффициентами}{\No\,10}


\tl{Баутин С.П., Казаков А.Л.}


{Одна задача Коши с начальными данными 


на разных \newline поверхностях для системы с особенностью}{\No\,10}


\tl{Бенкафадар Н.М., Гельман Б.Д.}{О локальной степени для одного


класса\newline


многозначных полей}{\No\,2}


\tl{Борисович А.Ю.}{Функционально-операторный метод исследования бифуркаций


в\newline


эквивариантной проблеме Плато}{\No\,2}


\tl{Борисович Ю.Г.}{Топологические характеристики и исследование


разрешимости\newline


нелинейных проблем}{\No\,2}


\tl{Бояршинова А.В.}


{О пространстве существенных инфинитезимальных деформаций}{\No\,8}


\tl{Бродский Г.М.}{Условие Гротендика $AB5^*$ и обобщения дистрибутивности 


модуля}{\No\,3}


\tl{Булатов М.В.}


{Метод возмущения дифференциально-алгебраических систем}{\No\,11}


\tl{Васильев А.Ф., Васильева Т.И.}


{О конечных группах, у которых главные факторы\newline


являются простыми группами}{\No\,11}


\tl{Васильев А.Ю.}{Изопериметрические теоремы роста для квазиконформных 


автомор-\newline физмов круга}{\No\,3}


\tl{Васильева Т.И., Васильев А.Ф.}


{О конечных группах, у которых главные факторы\newline


являются простыми группами}{\No\,11}


\tl{Велимирович Л.С.}{Бесконечно малые изгибания торообразной


поверхности вращения \newline


с многоугольным меридианом}{\No\,9}


\tl{Витюк А.Н.}{Существование решений дифференциальных включений 


с частными\newline производными дробных порядков}{\No\,8}


\tl{Володин А.И.}{Сходимость взвешенных сумм случайных элементов 


в банаховых\newline пространствах типа $p$}{\No\,8}


\tl{Ву Дж., Кравцевич В.}{Бифуркация систем, обратимых по времени}{\No\,2}


\tl{Вуколова Т.М., Дьяченко М.И.}{Оценки смешанных норм сумм двойных 


тригонометри-\newline


ческих рядов с кратно монотонными коэффициентами}{\No\,7}


\tl{Вьялицин А.А.}{Исследование систем линейных диофантовых


неравенств методами\newline теории групп}{\No\,11}


\tl{Габдулхаев Б.Г., Ермолаева Л.Б.}


{Интерполяционныe полиномы Лагранжа в прост-\newline


ранствах Соболева}{\No\,5}


\tl{Гейт В.Э.}{Об аппроксимационных свойствах 


высших производных периодических\newline функций}{\No\,10}


\tl{Гельман Б.Д., Бенкафадар Н.М.}{О локальной степени для одного


класса\newline


многозначных полей}{\No\,2}


\tl{Гимади Э.Х., Залюбовский В.В.}{Задача упаковки в контейнеры:


асимптотически\newline точный подход}{\No\,12} 


\tl{Гимади Э.Х., Залюбовский В.В., Шарыгин П.И.}{Задача упаковки в


полосу:\newline


асимптотически точный подход}{\No\,12}


\tl{Глазырина Л.Л., Павлова М.Ф.}{О разрешимости одного нелинейного


эволюционного\newline


неравенства теории совместного движения поверхностных и подземных вод}


{\No\,4}


\tl{Глушак А.В.}{Об асимптотической близости решений 


задачи Коши для дифференциаль-\newline


ных уравнений первого порядка в банаховом пространстве}{\No\,7}


\tl{Голиков А.А., Евтушенко Ю.Г.}{Модифицированная функция Лагранжа 


для задач\newline


линейного программирования}{\No\,12}


\tl{Горюнов Ю.Ю.}{Обобщение теоремы Дресслера}{\No\,7}


\tl{Данилов Ю.М., Никонова Н.В.} 


{Неявная разностная схема для решения нелинейных\newline


уравнений гидродинамики}{\No\,5}


\tl{Даньшин А.Ю.}{Инфинитезимальные проективные преобразования


в касательном\newline


расслоении общего пространства путей}{\No\,9}


\tl{Даринский Б.М., Сапронов Ю.И.}{Бифуркации экстремалей вблизи


особенности\newline


многомерной сборки}{\No\,2}


\tl{Дербенева О.В.}{Группы с относительно большой мерой коммутативности}


{\No\,3}


\tl{Джарадин Джехад.}{О формациях с системами наследственных подформаций}


{\No\,1}


\tl{Додунова Л.К.}{Приближение аналитических функций суммами Валле-Пуссена}


{\No\,3}


\tl{Дудко Л.Л., Свиридюк Г.А., Сукачева Т.Г.} 


{Относительная $\sigma$-ограниченность линей-\newline


ных операторов}{\No\,7}


\tl{Дьяконов Е.Г.}{Оценки $N$-поперечников в смысле Колмогорова для некоторых 


ком-\newline


пактов в усиленных пространствах Соболева}{\No\,4}


\tl{Дьяченко М.И., Вуколова Т.М.} 


{Оценки смешанных норм сумм двойных 


тригонометри-\newline


ческих рядов с кратно монотонными коэффициентами}{\No\,7}


\tl{Евтушенко Ю.Г., Голиков А.А.}{Модифицированная функция Лагранжа 


для задач\newline линейного программирования}{\No\,12}


\tl{Егишянц С.А., Островский Е.И.} 


{Асимптотическое поведение случайных полей 


на\newline последовательности расширяющихся множеств}{\No\,8}


\tl{Еремин И.И.}{Некоторые


вопросы кусочно-линейного программирования}{\No\,12}


\tl{Ермолаев Ю.Б.}{Универсальное картановское продолжение}{\No\,11}


\tl{Ермолаева Л.Б., Габдулхаев Б.Г.}


{Интерполяционныe полиномы Лагранжа в прост-\newline


ранствах Соболева}{\No\,5}


\tl{Ефремова Л.С., Махрова Е.Н.}


{О динамике монотонного отображения $n$-ода}{\No\,10}


\tl{Жевнова Н.Г., Скиба А.Н.} 


{$p$-насыщенные формации с дополняемыми $p$-насыщенными\newline


подформациями}{\No\,5}


\tl{Жогова Т.Б.}


{Проективное изгибание второго порядка семейств $L^m_{2n-1}$}{\No\,9}


\tl{Заботин И.Я.}{Алгоритм второго порядка с параметризованными


направлениями для\newline задач условной оптимизации}{\No\,12}


\tl{Забудский Г.Г.}{Алгоритм решения одной задачи


оптимального линейного упорядочения}{\No\,12}


\tl{Зайцева С.Б., Злотник А.А.}{Точные оценки градиента погрешности 


локально-одно-\newline мерных 


методов для многомерного уравнения теплопроводности}{\No\,4}


\tl{Залюбовский В.В., Гимади Э.Х.}{Задача упаковки в контейнеры:


асимптотически\newline точный подход}{\No\,12} 


\tl{Залюбовский В.В., Гимади Э.Х., Шарыгин П.И.}{Задача упаковки в


полосу:\newline


асимптотически точный подход}{\No\,12}


\tl{Звягин В.Г.}{Индекс нулевой точки вполне непрерывного возмущения 


фредгольмова\newline


отображения, коммутирующего с действием тора}{\No\,2}


\tl{Злотник А.А. Амосов А.А.}{Полудискретный метод решения уравнений 


одномерного\newline


движения неоднородного вязкого теплопроводного газа с негладкими данными}


{\No\,4}


\tl{Злотник А.А., Зайцева С.Б.}{Точные оценки градиента погрешности 


локально-одно-\newline мерных методов для многомерного уравнения теплопроводности}


{\No\,4}


\tl{Казаков А.Л., Баутин С.П.}


{Одна задача Коши с начальными данными 


на разных \newline поверхностях для системы с особенностью}{\No\,10}


\tl{Карцатос А.Г.}{Об условиях седловой точки в теории возмущений нелинейных 


макси-\newline мально монотонных операторов в гильбертовом пространстве}


{\No\,2}


\tl{Карчевский М.М., Ляшко А.Д., Паймушин В.Н.}{О математических задачах


теории\newline


многослойных оболочек с трансверсально-мягкими заполнителями}{\No\,4}


\tl{Кац М.М.}{Критерий линейной упорядочиваемости 


частичного автомата}{\No\,10}


\tl{Келлер А.В., Свиридюк Г.А.} 


{Инвариантные пространства и дихотомии решений одного\newline


класса линейных уравнений типа Соболева}{\No\,5}


\tl{Кельзон А.А.}


{О тригонометрическом интерполировании функций


$m$-гармонической \newline ограниченной вариации}{\No\,10}


\tl{Климов В.С.}{Теоремы вложения и внешние поперечные меры}{\No\,7}


\tl{Коннов И.В.}{О системах вариационных неравенств}{\No\,12}


\tl{Копаев А.В.}{Об одном интегральном преобразовании гармонических в 


полупространстве\newline функций}{\No\,1}


\tl{Косов А.А.}{О глобальной устойчивости неавтономных


систем. {\rm I}}{\No\,7}


\tl{Коcов А.А.}{О глобальной ycтойчивоcти


неавтономныx cиcтем. {\rm II}}{\No\,8}


\tl{Коцур И.М., Коцур М.Ф.}


{О радиусах $n$-листной выпуклости и звездности 


для \newline некоторых специальных классов аналитических 


функций в круге}{\No\,10}


\tl{Коцур М.Ф., Коцур И.М.}


{О радиусах $n$-листной выпуклости и звездности 


для \newline некоторых специальных классов аналитических 


функций в круге}{\No\,10}


\tl{Кравцевич В., Ву Дж.}{Бифуркация систем, обратимых по времени}{\No\,2}


\tl{Крукиер Л.А.}{Кососимметричные итерационные методы решения стационарной 


задачи\newline конвекции-диффузии с малым параметром при старшей производной}


{\No\,4}


\tl{Кунаковская О.В.}{О гладких разбиениях единицы 


на банаховых многообразиях}{\No\,10}


\tl{Лактионов С.А.}{О комплексах двумерных плоскостей в проективном


пространстве\newline


с одной критической прямой}{\No\,5}


\tl{Леденева Т.М.}


{Некоторые аспекты представления нечетких


операторов отношением\newline двух многочленов}{\No\,11}


\tl{Лернер А.К.}{Об оценках сильных максимальных функций}{\No\,7}


\tl{Липагина Л.В.}{О строении алгебры инвариантных


аффинорных структур на сфере $S^5$}{\No\,9}


\tl{Липачева Е.В.}{Частичные автоморфизмы счетных моделей}{\No\,1}


\tl{Лихачева Н.Н.}{Теорема Валле-Пуссена для одного класса 


функционально-дифференциаль-\newline


ных уравнений}{\No\,5}


\tl{Лосев А.Г.} 


{О взаимосвязи некоторых лиувиллевых теорем 


на римановых многообразиях \newline специального вида}{\No\,10}


\tl{Ляшко А.Д., Карчевский М.М., Паймушин В.Н.}{О математических задачах


теории\newline


многослойных оболочек с трансверсально-мягкими заполнителями}{\No\,4}


\tl{Майорова М.Е., Павлова М.Ф.}{Сходимость явных разностных схем 


для одного вариа-\newline


ционного неравенства 


теории нелинейной нестационарной фильтрации}{\No\,7}


\tl{Марзантович В.}{Нахождение периодических точек гладких отображений с 


использо-\newline ванием лефшецевых чисел и их степеней}{\No\,2}


\tl{Марков П.Е.}{Общие аналитические и


бесконечно малые деформации погружений. {\rm I}}{\No\,9}


\tl{Марков П.Е.}{Общие аналитические и бесконечно малые


деформации погружений. {\rm II}}{\No\,11}


\tl{Махнев А.А.}{Характеризация одного класса реберно-регулярных


графов}{\No\,1}


\tl{Махрова Е.Н., Ефремова Л.С.}


{О динамике монотонного отображения $n$-ода}{\No\,10}


\tl{Мелихов С.Н., Момм З.}{О линейном непрерывном правом обратном 


для оператора\newline


свертки на пространствах ростков 


аналитических функций на выпуклых компактах в $\Bbb C$}{\No\,5}


\tl{Мирзоян В.А.}{Подмногообразия с симметрическими


фундаментальными формами\newline


высшего порядка как огибающие}{\No\,9}


\tl{Митрохин С.И.}{О ``расщеплении'' кратных в главном собственных значений 


много-\newline точечных краевых задач}{\No\,3}


\tl{Михайличенко Г.Г.}{Трехмерные алгебры Ли локально транзитивных


преобразований\newline


пространства}{\No\,9}


\tl{Момм З., Мелихов С.Н.}{О линейном непрерывном правом обратном 


для оператора\newline


свертки на пространствах ростков 


аналитических функций на выпуклых компактах в $\Bbb C$}{\No\,5}


\tl{Непомнящих Ю.В., Шрагин И.В.} 


{$D$-условия Каратеодори и их связь 


с $D$-непрерыв-\newline


ностью оператора Немыцкого}{\No\,6}


\tl{Никонова Н.В., Данилов Ю.М.} 


{Неявная разностная схема для решения нелинейных\newline


уравнений гидродинамики}{\No\,5}


\tl{Новак С.Ю.}{Статистическое оценивание максимального 


собственного числа матрицы}{\No\,5}


\tl{Онегов Л.А.}{Об одной задаче оптимального восстановления интегралов в


смысле\newline


главного значения по Коши и конечного значения по Адамару}{\No\,1}


\tl{Орлов В.П.}{Слабо вырождающиеся дифференциальные уравнения с


неограниченным\newline


операторным коэффициентом}{\No\,1}


\tl{Орлов В.П.}{Коэрцитивная разрешимость слабо вырождающихся


дифференциальных\newline


уравнений с неограниченным операторным коэффициентом}{\No\,3}


\tl{Островский Е.И., Егишянц С.А.} 


{Асимптотическое поведение случайных полей 


на\newline последовательности расширяющихся множеств}{\No\,8}


\tl{Павлова М.Ф., Глазырина Л.Л.}{О разрешимости одного нелинейного


эволюционного\newline


неравенства теории совместного движения поверхностных и подземных вод}


{\No\,4}


\tl{Павлова М.Ф., Майорова М.Е.}{Сходимость явных разностных схем 


для одного вариа-\newline


ционного неравенства 


теории нелинейной нестационарной фильтрации}{\No\,7}


\tl{Паймушин В.Н., Карчевский М.М., Ляшко А.Д.}{О математических задачах


теории\newline многослойных оболочек с трансверсально-мягкими заполнителями}


{\No\,4}


\tl{Пайсон О.Б.}{Многообразия ассоциативных колец, все критические кольца


которых\newline являются арифметическими}{\No\,1}


\tl{Пейсахович Я., Фицпатрик П.М.}{Дополнение отображения в сингулярной


точке\newline


максимального ранга}{\No\,2}


\tl{Петрова В.В., Тонков Е.Л.}


{Допустимость периодических процессов и теоремы\newline


существования периодических решений. {\rm II}}{\No\,6}


\tl{Пинус А.Г.}{Об $\aleph$-идентичности отношений эпиморфности


и вложимости на квазимного-\newline образиях}{\No\,11}


\tl{Ратинер Н.М.} 


{Свойства классов бордизмов оснащенных


подмногообразий фредгольмовой\newline структуры}{\No\,8}


\tl{Репницкий В.Б.}


{О решеточно универсальных многообразиях алгебр}{\No\,5}


\tl{Рзаев Р.М.}{Многомерный сингулярный интегральный оператор в


пространствах\newline


$BMO_{\varphi,\theta}$ и $H_{\varphi,\theta}$}{\No\,3}


\tl{Румянцев А.Н.} 


{Конструктивное исследование дифференциальных уравнений\newline


с произвольным отклонением аргумента}{\No\,6}


\tl{Рязанцева И.П.} 


{Об уравнениях с возмущенными аккретивными отображениями}{\No\,7}


\tl{Сабитов И.Х.}


{Интегральный аналог уравнений Дарбу


для погружений двумерных метрик}{\No\,9}


\tl{Сапронов Ю.И., Даринский Б.М.}{Бифуркации экстремалей вблизи


особенности\newline


многомерной сборки}{\No\,2}


\tl{Свиридюк Г.А., Келлер А.В.} 


{Инвариантные пространства и дихотомии решений одного\newline


класса линейных уравнений типа Соболева}{\No\,5}


\tl{Свиридюк Г.А., Дудко Л.Л., Сукачева Т.Г.} 


{Относительная $\sigma$-ограниченность линей-\newline


ных операторов}{\No\,7}


\tl{Севастьянов В.А.}{Метод Римана для трехмерного гиперболического


уравнения третьего\newline


порядка}{\No\,5}


\tl{Седова С.М.}{О критерии устойчивости одного скалярного уравнения


с постоянным\newline запаздыванием и периодическим коэффициентом}{\No\,11}


\tl{Секерин А.Б.}{Теорема о носителе для комплексного 


преобразования Радона}{\No\,8}


\tl{Серовайский С.Я.}{Расширенно дифференцируемые подмногообразия}{\No\,1}


\tl{Сидоров И.Н.}{Построение интегрального представления 


решения уравнений равновесия\newline


различных вариантов теории оболочек сложной геометрии}{\No\,7}


\tl{Сильченко Ю.Т.}{Оператор дифференцирования 


с нерегулярными граничными условиями}{\No\,6}


\tl{Симонов П.М., Азбелев Н.В.}


{Устойчивость уравнений с запаздывающим аргументом}{\No\,6}


\tl{Симонов П.М., Чистяков А.В.} 


{Об экспоненциальной устойчивости линейных дифферен-\newline


циально-разностных систем}{\No\,6}


\tl{Скиба А.Н., Жевнова Н.Г.} 


{$p$-насыщенные формации с дополняемыми $p$-насыщенными\newline


подформациями}{\No\,5}


\tl{Сосов Е.Н.}{О геодезических отображениях 


специальных метрических пространств}{\No\,8}


\tl{Сукачева Т.Г., Дудко Л.Л., Свиридюк Г.А.} 


{Относительная $\sigma$-ограниченность линей-\newline


ных операторов}{\No\,7}


\tl{Сухочева Л.И.}


{Новый критерий принадлежности оператора классу $K(H)$}{\No\,10}


\tl{Терегулов И.Г.}


{Термодинамические потенциалы и математическое


моделирование\newline процесса течения вязких сред}{\No\,11}


\tl{Тонков Е.Л., Петрова В.В.}


{Допустимость периодических процессов и теоремы\newline


существования периодических решений. {\rm II}}{\No\,6}


\tl{Тонконог С.Л.}


{Об уравнениях динамики неньютоновской жидкости, 


движущейся с \newline проскальзыванием относительно ложа}{\No\,10}


\tl{Тронин С.Н.} 


{О Морита-контекстах, возникающих при изоморфизме


некоторых подколец \newline колец эндоморфизмов}{\No\,8}


\tl{Ульянова Е.Л.}


{О спектральных свойствах относительно 


конечномерных возмущений \newline самосопряженных операторов}{\No\,10}


\tl{Умалатов С.Д.} 


{Об одном пространстве К\"ете, 


связанном с дифференциальным оператором}{\No\,6}


\tl{Ускова Н.Б.,}{К одной проблеме Пирси и Шилдса}{\No\,10}


\tl{Фарафонова Н.К.}


{Геометрия грассманова расслоения. {\rm I}}{\No\,9}


\tl{Федотов Е.М.}{Об одном классе двухслойных разностных схем для нелинейных 


краевых\newline задач с памятью}{\No\,4}


\tl{Фицпатрик П.М., Пейсахович Я.}{Дополнение отображения в сингулярной точке 


макси-\newline мального ранга}{\No\,2}


\tl{Хамисов О.В., Апекина Е.В.}{Модифицированный метод симплексных


погружений\newline


с одновременным введением нескольких секущих плоскостей}{\No\,12}


\tl{Целищева И.В., Шишкин Г.И.}{Метод декомпозиции для сингулярно


возмущенных\newline


краевых задач 


с локальным возмущением начальных условий. Уравнения с


конвек-\newline тивными членами}


{\No\,4}


\tl{Ченцов А.Г.}{К вопросу о расширении стохастических


ограничений в классе конечно-\newline


аддитивных вероятностей}{\No\,6}


\tl{Чешкова М.А.}{О торсообразующем векторном поле}{\No\,1}


\tl{Чешкова М.А.}


{Преобразование Бианки $n$-поверхностей в $E^{2n-1}$}{\No\,9}


\tl{Чистяков А.В., Симонов П.М.} 


{Об экспоненциальной устойчивости линейных дифферен-\newline


циально-разностных систем}{\No\,6}


\tl{Шарыгин П.И., Гимади Э.Х., Залюбовский В.В.}{Задача упаковки в


полосу:\newline


асимптотически точный подход}{\No\,12}


\tl{Шевченко В.Н.}{О разбиении выпуклого политопа на симплексы


без новых вершин}{\No\,12}


\tl{Шишкин Г.И., Целищева И.В.}{Метод декомпозиции для сингулярно


возмущенных\newline


краевых задач с локальным возмущением начальных условий. 


Уравнения с конвек-\newline тивными членами}{\No\,4}


\tl{Шрагин И.В., Непомнящих Ю.В.}{$D$-условия Каратеодори и их связь 


с $D$-непрерыв-\newline


ностью оператора Немыцкого}{\No\,6}


\tl{Яковлев Е.И.}


{Секционные кривизны многообразий типа Калуцы-Клейна}{\No\,9}


\tl{Якупов М.Ш.}{Каноническое описание и обращение операторов 


в дискретном\newline пространстве. {\rm I}}{\No\,8}





\bigskip 


{\centerline {КРАТКИЕ СООБЩЕНИЯ}}


\bigskip


 


\tl{Абилов В.А., Абилова Ф.В.}{Приближение функций алгебраическими


многочленами\newline в среднем}{\No\,3}


\tl{Абилова Ф.В., Абилов В.А.}{Приближение функций алгебраическими


многочленами\newline в среднем}{\No\,3}


\tl{Авхадиев Ф.Г., Аксентьев Л.А., Бильченко Г.Г.}{Классы однолистных и 


много-\newline листных интегралов Кристоффеля-Шварца и их приложения}{\No\,3}


\tl{Агеев А.Л.}{Решение уравнений {\rm I}-рода с конечномерной нелинейностью}


{\No\,3}


\tl{Аксентьев Л.А., Авхадиев Ф.Г., Бильченко Г.Г.}{Классы однолистных и 


много-\newline листных интегралов Кристоффеля-Шварца и их приложения}


{\No\,3}


\tl{Аксентьев Л.А., Бильченко Г.Г.}{К обратной задаче для интегралов 


Кристоффеля--\newline Шварца}{\No\,8}


\tl{Бильченко Г.Г., Авхадиев Ф.Г., Аксентьев Л.А.}{Классы однолистных и 


много-\newline листных интегралов Кристоффеля-Шварца и их приложения}


{\No\,3}


\tl{Бильченко Г.Г., Аксентьев Л.А.} 


{К обратной задаче для интегралов 


Кристоффеля--\newline Шварца}{\No\,8}


\tl{Бильченко Н.Г., Овчинников В.А.}{О группе Ли, допускаемой одной


нелинейной\newline


системой уравнений в частных производных параболического типа}{\No\,1}


\tl{Бужуф Хамза, Мухин В.В.} 


{О топологиях на полугруппах и группах, определяемых\newline


семействами отклонений и норм}{\No\,5}


\tl{Гурьянов Н.Г.}{Двоякопериодическое решение задачи определения 


напряженно-дефор-\newline мированного состояния слоистой сферической оболочки}


{\No\,1}


\tl{Мартюшов С.Н.}{Построение дву- и трехмерных сеток для задач газодинамики на 


основе\newline уравнения Пуассона}{\No\,4}


\tl{Мухин В.В., Бужуф Хамза} 


{О топологиях на полугруппах и группах, определяемых\newline


семействами отклонений и норм}{\No\,5}


\tl{Овчинников В.А., Бильченко Н.Г.}{О группе Ли, допускаемой одной


нелинейной\newline


системой уравнений в частных производных параболического типа}{\No\,1}


\tl{Смирнова А.Б.}{Монотонные процессы для нелинейных некорректных задач в 


$K$-прост-\newline ранствах}{\No\,1}


\tl{Федотов А.И.}{Сходимость квадратурно-разностного метода для одного класса 


линейных\newline сингулярных интегро-дифференциальных уравнений на отрезке}


{\No\,3}


\bigskip





\noindent Рефераты статей, депонированных в ВИНИТИ через журнал ``Известия


вузов. Математика'' \dotfill \No\,1





\noindent Рефераты статей, депонированных в ВИНИТИ через журнал ``Известия


вузов. Математика'' \dotfill \No\,3





\end{document}


